\sporkSubsection{TMA and its Limitations}{sec:TmaLimitations}

The CUDA TMA (tensor memory accelerator) presents several unique code representation challenges and exposes limitations in our analysis.
The TMA instructions (\texttt{cp\{.reduce\}.async.bulk} \cite{ptx})
copy entire tiles between GMEM and SMEM in a single instruction execution.
Optionally, the SMEM tile is atomically reduced into GMEM.
Two features in the grammar (Section~\ref{sec:Grammar}, Figure~\ref{fig:stmt}) are unique to TMA usage.

First, TMA instructions do not take the GMEM pointer, the tile size, GMEM strides, or SMEM swizzle mode directly.
These attributes are packed as an opaque \texttt{CUtensorMap} on the CPU,
which is eventually passed to the TMA instruction.
We represent this with a ``special window statement'' in \sporkName{},
which constructs an alias to the underlying memory, in CPU scope.
The TMA instruction performs internal bound checking for the GMEM loads and stores:
out-of-bound loads fetch zeros or NaN (configurable), and out-of-bounds stores are masked out.

Second, TMA instructions in the GMEM-to-SMEM direction take a trailing barrier expression
referencing mbarriers (Section~\ref{sec:Mbarriers});
the instruction directly interacts with the mbarrier.
We chose to put the barrier separate from the ordinary parameters
to avoid modifying Exo's \texttt{replace} rule's unification,
in line with the view that this trailing barrier is an ``annotation'' on sequential code.

Our correctness checking of TMA usage is incomplete in multiple ways.
\textbf{(1)} We used unsafe rewrites to deal with Exo not being able to handle TMA's padding behavior
(Figure~\ref{fig:TmaUnsafeRewrite});
relatedly, we had to disable bounds-checking on a case-by-case basis to enable TMA usage.
\textbf{(2)} Internally, the \sporkName{} TMA instruction expands to the actual TMA instruction and a hidden expect-tx instruction, both taking the mbarriers as parameters.
We fail to decouple two instructions,
and do not reason that this mbarrier usage interacts correctly (is race free) with ordinary arrives and awaits of the mbarrier.
\textbf{(3)} The atomic reduction somewhat undermines the sequential-parallel equivalence argument (Section~\ref{sec:TraceAtomics}).

% Putting the entire discussion here to save horizontal space and
% because it's interrupting the main flow of the paper.

\begin{figure}[h!]
\sporkCodeRule
\begin{minipage}{0.275\linewidth}
\sporkListing{tma_mini_unsafe}
\end{minipage}
\begin{minipage}{0.700\linewidth}
\vspace{2pt}
The listing shows a minimal example of the sort of situtation that currently forces an unsafe rewrite for TMA usage.
The situation arises when the \texttt{divide\_loop} (split) rewrite creates an inner loop with a guard
to account for the tail case.
Even though removing the example's first guard is safe since \texttt{tmp[6]} and \texttt{tmp[7]} are not being read into \texttt{C}
in the second loop thanks to the second guard,
Exo's dependency-based checks treat these writes as semantically relevant and cannot prove they are observationally redundant.
A similar issue occurs when removing guards inside accumulator loops,
where accumulating ``past the end'' is safe due to TMA padding with zeros.
We marked these ``remove guard'' rewrites as \emph{unsafe} in our schedules.
Both issues would be resolved by future work extending Exo's sequential rewrites with value-sensitivity.
\end{minipage}
\caption{TMA unsafe rewrite minimal reproducer \& discussion}
\label{fig:TmaUnsafeRewrite}
\sporkCodeRule
\end{figure}

